\UseRawInputEncoding
\documentclass[a4paper, 14pt]{extarticle}

% Поля
%--------------------------------------
\usepackage{geometry}
\geometry{a4paper,tmargin=2cm,bmargin=2cm,lmargin=3cm,rmargin=1cm}
%--------------------------------------


%Russian-specific packages
%--------------------------------------
\usepackage[T2A]{fontenc}
\usepackage[utf8]{inputenc} 
\usepackage[english, main=russian]{babel}
%--------------------------------------

\usepackage{textcomp}

% Красная строка
%--------------------------------------
\usepackage{indentfirst}               
%--------------------------------------             


%Graphics
%--------------------------------------
\usepackage{graphicx}
\graphicspath{ {./images/} }
\usepackage{wrapfig}
%--------------------------------------

% Полуторный интервал
%--------------------------------------
\linespread{1.3}                    
%--------------------------------------

%Выравнивание и переносы
%--------------------------------------
% Избавляемся от переполнений
\sloppy
% Запрещаем разрыв страницы после первой строки абзаца
\clubpenalty=10000
% Запрещаем разрыв страницы после последней строки абзаца
\widowpenalty=10000
%--------------------------------------

%Списки
\usepackage{enumitem}

%Подписи
\usepackage{caption} 

%Гиперссылки
\usepackage{hyperref}

\hypersetup {
	unicode=true
}

%Рисунки
%--------------------------------------
\DeclareCaptionLabelSeparator*{emdash}{~--- }
\captionsetup[figure]{labelsep=emdash,font=onehalfspacing,position=bottom}
%--------------------------------------

\usepackage{tempora}
\usepackage{amsmath}
\usepackage{color}
\usepackage{listings}
\lstset{
  belowcaptionskip=1\baselineskip,
  breaklines=true,
  frame=L,
  xleftmargin=\parindent,
  language=Python,
  showstringspaces=false,
  basicstyle=\footnotesize\ttfamily,
  keywordstyle=\bfseries\color{blue},
  commentstyle=\itshape\color{purple},
  identifierstyle=\color{black},
  stringstyle=\color{red},
  extendedchars=\true,
}

%--------------------------------------
%			НАЧАЛО ДОКУМЕНТА
%--------------------------------------

\begin{document}

%--------------------------------------
%			ТИТУЛЬНЫЙ ЛИСТ
%--------------------------------------
\begin{titlepage}
\thispagestyle{empty}
\newpage


%Шапка титульного листа
%--------------------------------------
\vspace*{-60pt}
\hspace{-65pt}
\begin{minipage}{0.3\textwidth}
\hspace*{-20pt}\centering

\includegraphics[width=\textwidth]{emblem}
\end{minipage}
\begin{minipage}{0.67\textwidth}\small \textbf{
\vspace*{-0.7ex}
\hspace*{-6pt}\centerline{Министерство науки и высшего образования Российской Федерации}
\vspace*{-0.7ex}
\centerline{Федеральное государственное бюджетное образовательное учреждение }
\vspace*{-0.7ex}
\centerline{высшего образования}
\vspace*{-0.7ex}
\centerline{<<Московский государственный технический университет}
\vspace*{-0.7ex}
\centerline{имени Н.Э. Баумана}
\vspace*{-0.7ex}
\centerline{(национальный исследовательский университет)>>}
\vspace*{-0.7ex}
\centerline{(МГТУ им. Н.Э. Баумана)}}
\end{minipage}
%--------------------------------------

%Полосы
%--------------------------------------
\vspace{-25pt}
\hspace{-35pt}\rule{\textwidth}{2.3pt}

\vspace*{-20.3pt}
\hspace{-35pt}\rule{\textwidth}{0.4pt}
%--------------------------------------

\vspace{1.5ex}
\hspace{-35pt} \noindent \small ФАКУЛЬТЕТ\hspace{80pt} <<Информатика и системы управления>>

\vspace*{-16pt}
\hspace{47pt}\rule{0.83\textwidth}{0.4pt}

\vspace{0.5ex}
\hspace{-35pt} \noindent \small КАФЕДРА\hspace{50pt} <<Теоретическая информатика и компьютерные технологии>>

\vspace*{-16pt}
\hspace{30pt}\rule{0.866\textwidth}{0.4pt}
  
\vspace{11em}

\begin{center}
\Large {\bf Лабораторная работа № 5} \\ 
\large {\bf<<Вычисление собственных значений и собственных векторов симметричной матрицы методом А.М. Крылова>>} \\ 
\large {\bf по курсу <<Численные методы линейной алгебры>>} \\ 
\end{center}\normalsize

\vspace{8em}


\begin{flushright}
  {Студент группы ИУ9-71Б Лысенко В. В.\hspace*{15pt} \\
  \vspace{2ex}
  Преподаватель Посевин Д. П.\hspace*{15pt}}
\end{flushright}

\bigskip

\vfill
 

\begin{center}
\textsl{Москва 2024}
\end{center}
\end{titlepage}
%--------------------------------------
%		КОНЕЦ ТИТУЛЬНОГО ЛИСТА
%--------------------------------------

\renewcommand{\ttdefault}{pcr}

\setlength{\tabcolsep}{3pt}
\newpage
\setcounter{page}{2}

\section{Цель работы}\label{Sect::task}
\par
\renewcommand{\labelenumii}{\arabic{enumi}.\arabic{enumii}.}
Реализовать метод вычисления собственных значений и собственных векторов симметричной матрицы методом А.М. Крылова.

\section{Задачи}\label{Sect::task}
\par
\renewcommand{\labelenumii}{\arabic{enumi}.\arabic{enumii}.}
\begin{enumerate} 
    \item Реализовать метод вычисления собственных значений и собственных векторов симметричной матрицы методом А.М. Крылова.
    \item Реализовать метод поиска собственных значений действительной симметричной матрицы A размером 4х4.
    \item Проверить решение на матрице приведенной в презентации.
    \item Продемонстрировать работу приложения для произвольных симметричных матриц размером n х n с учетом выполнения пунктов приведенных выше.
    \item Сравнить результаты с полученными методом Данилевского.
\end{enumerate}

\pagebreak
\section{Практическая реализация}
\par
Сравнение результатов, полученных методами Данилевского и Крылова, реализовано по следущему принципу:
\par
При сравнении собственных значений вычисляются их разницы с собственными значениями, полученными библиотечным методом. После берётся самая большая разница для каждого из методов.
\par
При сравнении собственных векторов вычисляются углы между векторами, полученными каждым из методов, с векторами, вычесленными библиотечным методом. После берётся самый большой угол для каждого из методов. Углы вычеслены в радианнах.
\par
Ниже представлен код программы:
\begin{lstlisting}
using Random
using LinearAlgebra
using Plots
using Symbolics
using Polynomials

eps = 1e-10

function generate_symmetric_matrix(n)
    A = rand(Float64, n, n)
    for i in 1:n
        for j in i+1:n
            A[i,j] = A[j, i]
        end
    end  
    return A
end

function roots_for_Daniel(coefficients)
    x = -reverse(coefficients)
    append!(x, 1)
    p = Polynomial(x)
    
    return roots(p)
end


function danilevsky_method(A)
    n = size(A, 1)
    lambda_values = Vector{Float64}()
    lambda_vectors = Vector{Vector{Float64}}()
    main_B = I(n)

    B = A

    for k in 1:n-1
        B = zeros(Float64, n, n)
        for i in 1:n
            B[i,i] = 1.0
        end
        for j in 1:n
            B[n-k,j] = - A[n+1-k,j] / A[n+1-k, n-k]
            if j==n-k
                B[n-k,j] = 1 / A[n+1-k, n-k]
            end
        end

        A = inv(B) * A * B

        main_B *= B        
    end

    lambda_values = roots_for_Daniel(A[1, :])
    
    for val in lambda_values
        my_vector = Vector{Float64}()
        for i in 0:n-1
            push!(my_vector, val^i)
        end
        
        my_vector = main_B * reverse(my_vector)
        
        push!(lambda_vectors, my_vector / norm(my_vector))
    end
    return lambda_values, lambda_vectors
end

function roots_for_Krylov(coefficients)
    x = -coefficients
    append!(x, 1)
    p = Polynomial(x)
    
    return roots(p)
end


function krylov_method(A)
    n = size(A, 1)
    c = zeros(Float64, n)
    c[1] = 1.0
    c_vectors = Vector{Vector{Float64}}()
    push!(c_vectors, c)
    lambda_vectors = Vector{Vector{Float64}}()

    for i in 1:n
        c = A * c
        push!(c_vectors, c)
    end

    с_matrix = zeros(Float64, n, n)

    for i in 1:n
        for j in 1:n
            с_matrix[i,j] = c_vectors[j][i] 
        end
    end

    
    Alu = qr(с_matrix)
    p = Alu \ c_vectors[n+1]

    lambda_values = roots_for_Krylov(p)

    p = reverse(p)
    for val in lambda_values
        
        betas = Vector{Float64}()
        push!(betas, 1)
        for i in 2:n
            push!(betas, val*betas[i-1] - p[i-1]) 
            
        end

        
        x = reverse(с_matrix) * betas
        push!(lambda_vectors, reverse(x / norm(x) * sign(x[1])))
        
    end

    return lambda_values, lambda_vectors

end

function check_viete(lambda_values, A)
    trace_A = tr(A)
    det_A = det(A)
    sum_lambda_values = sum(lambda_values)
    product_lambda_values = prod(lambda_values)

    if isapprox(sum_lambda_values, trace_A) && isapprox(product_lambda_values, det_A)
        println("Проверка теоремы Виета пройдена")
    else
        println("Проверка теоремы Виета НЕ пройдена")
    end
end

function check_gershgorin(lambda_values, A)
    n = size(A, 1)
    is_valid = true

    for i in 1:n
        row_sum = sum(abs.(A[i, :])) - abs(A[i, i])
        if abs(lambda_values[i] - A[i, i]) > row_sum
            is_valid = false
        end
    end

    if is_valid
        println("Проверка кругов Гершгорина пройдена")
    else
        println("Проверка кругов Гершгорина НЕ пройдена")
    end
end


function are_orthogonal(vectors)
    n = size(vectors, 1)

    is_ortagonal = true
    for i in 1:n
        for j in 1:n
            if i != j
                dot_product = dot(vectors[i], vectors[j])
                if eps > dot_product
                    dot_product = 0
                end
                is_ortagonal = isapprox(dot_product, 0)
            end
            if !is_ortagonal
                break
            end
        end
        if !is_ortagonal
            break   
        end
    end

    if is_ortagonal
        println("Собственные вектора ортогональны")
    else
        println("Собственные вектора не ортогональны")
    end
end

function compare_values(lambda_values_danilevsky, lambda_values_krylov)
    diff = abs.(lambda_values_danilevsky - lambda_values_krylov)
    return maximum(diff)
end

function compare_vectors(lambda_vectors_danilevsky, lambda_vectors_krylov, n)
    angels = Vector{Float64}()
    for i in 1:n
        push!(angels, dot(lambda_vectors_danilevsky[i], lambda_vectors_krylov[i]) / (norm(lambda_vectors_danilevsky[i]) * norm(lambda_vectors_krylov[i]))
    end
    return maximum(angels)
end

A = [2.2 1 0.5 2;
     1 1.3 2 1;
     0.5 2 0.5 1.6;
     2 1 1.6 2]
A = Float64.(copy(A))
println("Симметричная матрица A:")
for i in 1:4
    for j in 1:4
        print(A[i,j], "  ")
    end
    println()
end
println()

lambda_values_A_Krylov, lambda_vectors_A_Krylov = krylov_method(A)
println("Собственные значения по Крылову:")
println(lambda_values_A_Krylov)
println()
println("Собственные векторы по Крылову:")
for vec in lambda_vectors_A_Krylov
    println(vec)
end
println()


check_viete(lambda_values_A_Krylov, A)
check_gershgorin(lambda_values_A_Krylov, A)
are_orthogonal(lambda_vectors_A_Krylov)
println()

println("-------------------------------------------------------------------------------")

lambda_values_A_Daniel, lambda_vectors_A_Daniel = danilevsky_method(A)
println("Собственные значения по Данилевскому:")
println(lambda_values_A_Daniel)
println()
println("Собственные векторы по Данилевскому:")
for vec in lambda_vectors_A_Daniel
    println(vec)
end
println()

check_viete(lambda_values_A_Daniel, A)
check_gershgorin(lambda_values_A_Daniel, A)
are_orthogonal(lambda_vectors_A_Daniel)
println()

compare_values(lambda_values_A_Daniel, lambda_values_A_Krylov)

println("-------------------------------------------------------------------------------")
println("-------------------------------------------------------------------------------")

B_size = 6
B = generate_symmetric_matrix(B_size)
println("Симметричная матрица B:")
for i in 1:B_size
    for j in 1:B_size
        print(B[i,j], "  ")
    end
    println()
end
println()

lambda_values, lambda_vectors = krylov_method(B)
println("Собственные значения по Крылову:")
println(lambda_values)
println()
println("Собственные векторы по Крылову:")
for vec in lambda_vectors
    println(vec)
end
println()

check_viete(lambda_values, B)
check_gershgorin(lambda_values, B)
are_orthogonal(lambda_vectors)
println()

check_viete(lambda_values, B)
check_gershgorin(lambda_values, B)
are_orthogonal(lambda_vectors)
println()

println("-------------------------------------------------------------------------------")

lambda_values, lambda_vectors = danilevsky_method(B)
println("Собственные значения по Данилевскому:")
println(lambda_values)
println()
println("Собственные векторы по Данилевскому:")
for vec in lambda_vectors
    println(vec)
end
println()

check_viete(lambda_values, B)
check_gershgorin(lambda_values, B)
are_orthogonal(lambda_vectors)
println()


cup = 15

Daniel_values_diffs = Vector{Float64}()
Daniel_vectors_diffs = Vector{Float64}()
Krylov_values_diffs = Vector{Float64}()
Krylov_vectors_diffs = Vector{Float64}()
cups = Vector{Float64}()

for n in 3:cup
    push!(cups, n)
    Array = generate_symmetric_matrix(n)
    
    classic = eigen(Array)
    classic_values = classic.values
    classic_vectors = Vector{Vector{Float64}}()
    for i in 1:n
        vec = Vector{Float64}()
        for j in 1:n
            push!(vec, classic.vectors[j,i])
        end
        push!(classic_vectors, vec)
    end

    
    Krylov_values, Krylov_vectors = krylov_method(Array)
    Daniel_values, Daniel_vectors = danilevsky_method(Array)

    push!(Krylov_values_diffs, compare_values(classic_values, Krylov_values))
    push!(Daniel_values_diffs, compare_values(classic_values, Daniel_values))
    push!(Krylov_vectors_diffs, compare_vectors(classic_vectors, Krylov_vectors, n))
    push!(Daniel_vectors_diffs, compare_vectors(classic_vectors, Daniel_vectors, n))
end


println()

p1 = plot(cups, Krylov_values_diffs, label="Krylov values", xlabel="Matrix size", ylabel="max error")
plot!(cups, Daniel_values_diffs, label="Daniel values", color=:red)
xticks!(0:1:15)
yticks!(0:0.01:1)
display(p1)

p3 = plot(cups, Krylov_vectors_diffs, label="Krylov vectors", xlabel="Matrix size", ylabel="angle error")
plot!(cups, Daniel_vectors_diffs, label="Daniel vectors", color=:red)
xticks!(0:1:15)
display(p3)
\end{lstlisting}

\section{Результаты}
\par
\begin{center}
\includegraphics[scale=0.65]{g1}
\end{center}
\par
\begin{center}
\includegraphics[scale=0.45]{g2}
\end{center}
\par
\begin{center}
\includegraphics[scale=0.65]{g3}
\end{center}
\par

\section{Выводы}
В ходе выполнения лаборатороной работы был реализован метод А.М. Крылова по вычислению собсвтвенных значений и собственных векторов. Метод был реализован для симметричных матриц произвольного размера. Корректность всех вычислений была проверена на матрице, приведённой в презентации. Также результаты вычеслений были сравнены с результатами метода Данилевского. Из приведённых графиков можно заключить, что метод Данилевского более точно вычисляет собственные значения матрицы, в то время как метод Крылова более точен в вычислении собственных векторов.
\end{document}
